\subsection{THE UNIFORMITY OF UNIFORM CONVERGENCE}

Let X be a set and let Y be a uniform space. For each entourage V of Y, let \(W(V)\) denote the set of all pairs \((u, v)\) of mappings of X into Y such that \((u(x), v(x)) \in V\) for all \(x \in X\) . As V runs through the set of entourages of Y, the sets \(W(V)\) form a fundamental system of entourages of a uniformity on F(X; Y). For they clearly satisfy Axiom \((\mathrm{U}_{\mathrm{I}}^{\prime})\) (Chapter II, § 1, no. 1); if V, \(V'\) are two entourages of Y such that \(V \subset V'\) , we have \(W(V) \subset W(V')\) , and therefore the sets \(W(V)\) satisfy \((\mathrm{B}_{\mathrm{I}})\) (Chapter I, § 6, no. 3); we have \(\widehat{W(V)} = W(\overline{V}^{-1})\) so that \((\mathrm{U}_{\mathrm{II}}^{\prime})\) is satisfied; finally, the relations `` \((u(x), v(x)) \in V\) for all \(x \in X\)'' and `` \((v(x), w(x)) \in V\) for all \(x \in X\)'' imply the relation `` \((u(x), w(x)) \in V^{2}\) for all \(x \in X\)'' ; in other words, we have \(\widehat{W(V)} \subset W(\overline{V}^{2})\) , which proves \((\mathrm{U}_{\mathrm{III}}^{\prime})\) .

DEFINITION 1. The uniformity on the set \(\mathcal{F}(\mathbf{X};\mathbf{Y})\) which has as a fundamental system of entourages the set of subsets \(W(\mathbf{V})\) , where \(\mathbf{V}\) runs through the set of entourages of \(\mathbf{Y}\) , is called the uniformity of uniform convergence. The topology it induces is called the topology of uniform convergence. If a filter \(\Phi\) on \(\mathcal{F}(\mathbf{X};\mathbf{Y})\) converges to an element \(u_0\) with respect to this topology, \(\Phi\) is said to converge uniformly to \(u_0\) .

Note that the topology of uniform convergence on \(\mathcal{F}(\mathrm{X};\mathrm{Y})\) depends on the uniform structure of Y and not merely on the topology of Y (Exercise 4).

The uniform space obtained by endowing \(\mathcal{F}(\mathbf{X};\mathbf{Y})\) with the uniformity of uniform convergence is denoted by \(\mathcal{F}_u(\mathbf{X};\mathbf{Y})\) .

